% ============================================================
%  MODELO DE PLANO DE NEGOCIOS -- Padrão ABNT (abntex2)
%  Baseado em: Dornelas, Sebrae, Pirola
%  Uso livre por qualquer empreendedor
% ============================================================

\documentclass[
	12pt,
	openright,
	oneside,
	a4paper,
	chapter=TITLE,
	section=TITLE,
	english,
	brazil
	]{abntex2}

% --- Pacotes essenciais ---
\usepackage{newtxtext,newtxmath}          % Times New Roman (ABNT)
\usepackage[T1]{fontenc}
\usepackage[utf8]{inputenc}
\usepackage{indentfirst}
\usepackage{xcolor}                       % Cores (substitui color)
\usepackage{colortbl}                     % Cores em tabelas
\usepackage{graphicx}
\usepackage{microtype}
\usepackage{booktabs}
\usepackage{tabularx}
\usepackage{multirow}
\usepackage{float}
\usepackage{amsmath}
\usepackage{enumitem}
\usepackage{setspace}                     % Controle de espaçamento
\usepackage{csquotes}                     % Aspas sensíveis ao idioma
\usepackage[style=abnt, backend=biber]{biblatex}
\usepackage{pgfplots}
\pgfplotsset{compat=1.18}

% --- Arquivo de referências ---
\addbibresource{referencias.bib}

% --- Definição de cores ---
\definecolor{azul-abnt}{RGB}{41,5,195}
\definecolor{azul-claro}{RGB}{235,238,255}   % Fundo de tabelas
\definecolor{cinza-instrucao}{gray}{0.50}
\definecolor{cinza-borda}{gray}{0.70}

% ============================================================
%  COMANDOS CUSTOMIZADOS
% ============================================================

% Campo preenchível com largura fixa (uso inline)
\newcommand{\campo}[2][6cm]{%
	\underline{\makebox[#1][l]{#2}}%
}

% Campo preenchível que ocupa o restante da linha (listas)
\newcommand{\icampo}{%
	\hrulefill
}

% Campo preenchível para células de tabela (largura flexível)
\newcommand{\tcampo}[1]{%
	{#1}%
}

% Instrução formatada com caixa sutil
\newcommand{\instrucao}[1]{%
	\vspace{0.3cm}%
	\noindent%
	{%
		\setlength{\fboxsep}{6pt}%
		\colorbox{azul-claro}{%
			\begin{minipage}{\dimexpr\textwidth-2\fboxsep\relax}%
				\setlength{\parindent}{1.25cm}%
				{\color{cinza-instrucao}\small\itshape #1}%
			\end{minipage}%
		}%
	}%
	\vspace{0.3cm}%
}

% Linha de preenchimento (para formulários)
\newcommand{\linha}[1][0.4pt]{%
	\rule{\linewidth}{#1}%
}

% ============================================================
%  CONFIGURAÇÕES ABNT
% ============================================================

% Margens ABNT (confirmadas)
% abntex2 já define: esquerda=3cm, direita=2cm, topo=3cm, base=2cm

% Espaçamento 1.5 entrelinhas (padrão ABNT)
\OnehalfSpacing

% Parágrafo: recuo de 1,25 cm (ABNT)
\setlength{\parindent}{1.25cm}
\setlength{\parskip}{0.0cm}   % sem espaço extra entre parágrafos

% Suprimir avisos de overfull/underfull (campos preenchíveis produzem linhas vazias)
\hbadness=10000
\hfuzz=1pt

% ============================================================
%  ESTILO DE CABEÇALHO / RODAPÉ (memoir, sem fancyhdr)
% ============================================================

\makepagestyle{planoheader}
\makeoddhead{planoheader}%
	{\small\itshape Plano de Negócios}%
	{}%
	{\small\thepage}
\makeheadrule{planoheader}{\textwidth}{0.4pt}
\makefootrule{planoheader}{\textwidth}{0pt}{}

\newcommand{\configurarcabecalho}{%
	\pagestyle{planoheader}%
}

% ============================================================
%  ESTILO DE CAPÍTULOS E SEÇÕES (ABNT)
% ============================================================

% Formato dos títulos de capítulo: centralizado, negrito, caixa alta
% O pacote abntex2 já aplica boa parte. Reforçamos:
\renewcommand{\ABNTEXchapterfont}{\normalfont\bfseries}
\renewcommand{\ABNTEXchapterfontsize}{\Large}

% ABNT: títulos de seção alinhados à esquerda, negrito, caixa alta
\renewcommand{\ABNTEXsectionfont}{\normalfont\bfseries}
\renewcommand{\ABNTEXsectionfontsize}{\large}

% ============================================================
%  INFORMAÇÕES DA CAPA / FOLHA DE ROSTO
% ============================================================

\titulo{Plano de Negócios}
\autor{[Nome do Empreendedor ou Equipe]}
\local{[Cidade -- UF]}
\data{[\today]}
\instituicao{%
	[Nome da Empresa]
	\par
	[Setor de Atuação / Segmento]%
}
\preambulo{Plano de negócios apresentado como ferramenta de planejamento e captação de recursos para o empreendimento [Nome da Empresa].}

% --- Configurações de hiperlinks e float ---
\makeatletter
\hypersetup{
	pdftitle={\@title},
	pdfauthor={\@author},
	pdfsubject={Plano de Negócios},
	pdfcreator={LaTeX with abnTeX2},
	colorlinks=true,
	linkcolor=azul-abnt,
	citecolor=azul-abnt,
	urlcolor=azul-abnt,
	bookmarksdepth=4
}
\setlength{\@fptop}{5pt}
\makeatother

% --- Quadros ABNT ---
\newcommand{\quadroname}{Quadro}
\newcommand{\listofquadrosname}{Lista de quadros}
\newfloat{quadro}{loq}{\quadroname}
\newlistof{listofquadros}{loq}{\listofquadrosname}
\newlistentry{quadro}{loq}{0}
\setfloatadjustment{quadro}{\centering}
\counterwithout{quadro}{chapter}
\renewcommand{\cftquadroname}{\quadroname\space}
\renewcommand*{\cftquadroaftersnum}{\hfill--\hfill}
\setfloatlocations{quadro}{hbtp}

% --- Índice ---
\makeindex

% ============================================================
%  INÍCIO DO DOCUMENTO
% ============================================================
\begin{document}

\selectlanguage{brazil}
\frenchspacing

% --- Capa / Folha de Rosto ---
\imprimirfolhaderosto

% --- Sumário ---
\tableofcontents*
\clearpage

% --- Configurar cabeçalho após capa ---
\configurarcabecalho

% ============================================================
%  ELEMENTOS TEXTUAIS
% ============================================================
\textual

% -------------------------------------------------------
%  1. SUMÁRIO EXECUTIVO
% -------------------------------------------------------
\chapter{Sumário Executivo}

\instrucao{Esta seção deve ser escrita por último. Resumo completo do plano em 1--2 páginas. Destaque os pontos mais importantes para atrair o leitor (investidor, sócio, banco).}

\section{Conceito do Negócio}

\centering
\begin{tabularx}{\textwidth}{r X}
\textbf{Nome da Empresa:} & \campo[8cm]{} \\[0.3cm]
\textbf{Setor:}           & \campo[5cm]{} \\[0.3cm]
\textbf{Porte:}           & \campo[5cm]{ME / EPP / MEI} \\[0.3cm]
\textbf{Forma Jurídica:}  & \campo[5cm]{LTDA / EIRELI / SLU} \\[0.3cm]
\textbf{Localização:}     & \campo[8cm]{} \\
\end{tabularx}

\vspace{0.5cm}
\textbf{Missão:} \campo[12cm]{}

\vspace{0.3cm}
\textbf{Visão:} \campo[12cm]{}

\vspace{0.3cm}
\textbf{Valores:} \campo[12cm]{}

\section{Oportunidade}

\instrucao{Qual o problema não resolvido do mercado? Como seu negócio resolve? Qual a janela de oportunidade?}

\section{Mercado e Vantagens Competitivas}

\instrucao{Tamanho do mercado, público-alvo, concorrentes e diferenciais.}

\section{Equipe de Gestão}

\instrucao{Principais sócios e executivos, formação e experiência.}

\section{Projeções Financeiras}

\vspace{0.2cm}
\centering
\begin{tabularx}{\textwidth}{X r}
\toprule
\rowcolor{azul-claro}
\textbf{Indicador} & \textbf{Valor} \\
\midrule
Investimento Inicial              & R\$ \hrulefill \\
Faturamento Estimado (Ano 1)     & R\$ \hrulefill \\
Margem Líquida Esperada           & \hrulefill\ \% \\
Payback Estimado                  & \hrulefill\ meses \\
TIR Estimada                      & \hrulefill\ \% a.a. \\
\bottomrule
\end{tabularx}
\label{tab:projecoes-financeiras}

\section{Aporte Necessário}

\instrucao{Recursos necessários, finalidade, prazo e contrapartida oferecida.}

\clearpage

% -------------------------------------------------------
%  2. CONCEITO DO NEGÓCIO
% -------------------------------------------------------
\chapter{Conceito do Negócio}

\section{Histórico da Empresa}

\instrucao{Como surgiu a ideia? Há empresa constituída ou é novo empreendimento?}

\section{Missão, Visão e Valores}

\centering
\begin{tabularx}{\textwidth}{r X}
\textbf{Missão:}   & \tcampo{} \\[0.3cm]
\textbf{Visão:}    & \tcampo{} \\[0.3cm]
\textbf{Valores:}  & \tcampo{} \\
\end{tabularx}

\section{Aspectos Legais e Societários}

\centering
\begin{tabularx}{\textwidth}{r X}
\textbf{Forma Jurídica:}       & \tcampo{} \\[0.3cm]
\textbf{Composição Societária:} & \tcampo{} \\[0.3cm]
\textbf{CNPJ:}                 & \tcampo{} \\[0.3cm]
\textbf{Registro / Licenças:}  & \tcampo{} \\
\end{tabularx}

\section{Localização e Abrangência}

\instrucao{Onde a empresa está ou se localizará? Qual área geográfica atenderá inicialmente?}

\section{Parcerias e Terceirizações}

\instrucao{Fornecedores estratégicos, parceiros de negócios, consultorias.}

\clearpage

% -------------------------------------------------------
%  3. ANÁLISE DE MERCADO
% -------------------------------------------------------
\chapter{Análise de Mercado}

\section{Análise Setorial (Macroambiente)}

\instrucao{Análise PEST: fatores Políticos, Econômicos, Sociais e Tecnológicos que afetam o setor.}

\centering
\begin{tabularx}{\textwidth}{X X}
\toprule
\rowcolor{azul-claro}
\textbf{Fator} & \textbf{Impacto no Negócio} \\
\midrule
Político-Legal  & \tcampo{} \\
Econômico       & \tcampo{} \\
Sociocultural   & \tcampo{} \\
Tecnológico     & \tcampo{} \\
\bottomrule
\end{tabularx}

\section{Mercado-Alvo}

\subsection{Segmentação de Clientes}

\instrucao{Quem compra, o que compra, por que compra, como compra e onde está?}

\centering
\begin{tabularx}{\textwidth}{X X}
\toprule
\rowcolor{azul-claro}
\textbf{Variável} & \textbf{Descrição} \\
\midrule
Perfil Demográfico  & \tcampo{} \\
Comportamento       & \tcampo{} \\
Necessidades        & \tcampo{} \\
Região Geográfica   & \tcampo{} \\
\bottomrule
\end{tabularx}

\subsection{Tamanho de Mercado}

\centering
\begin{tabularx}{\textwidth}{l X}
\toprule
\rowcolor{azul-claro}
\textbf{Tipo} & \textbf{Valor} \\
\midrule
Mercado Total (R\$)          & \tcampo{} \\
Mercado Potencial (R\$)      & \tcampo{} \\
Participação Pretendida      & \tcampo{} \\
Número Estimado de Clientes  & \tcampo{} \\
\bottomrule
\end{tabularx}
\label{tab:tamanho-mercado}

\section{Análise da Concorrência}

\instrucao{Concorrentes diretos e indiretos. Compare com seu negócio nos 4Ps.}

\centering
\begin{tabularx}{\textwidth}{X c c c c}
\toprule
\rowcolor{azul-claro}
\textbf{Concorrente} & \textbf{Qualidade} & \textbf{Preço} & \textbf{Local} & \textbf{Variedade} \\
\midrule
\tcampo{} & \tcampo{} & \tcampo{} & \tcampo{} & \tcampo{} \\
\tcampo{} & \tcampo{} & \tcampo{} & \tcampo{} & \tcampo{} \\
\tcampo{} & \tcampo{} & \tcampo{} & \tcampo{} & \tcampo{} \\
\textbf{Sua Empresa} & \tcampo{} & \tcampo{} & \tcampo{} & \tcampo{} \\
\bottomrule
\end{tabularx}
\label{tab:analise-concorrencia}

\section{Vantagens Competitivas}

\instrucao{O que torna seu negócio único? Por que o cliente escolherá você?}

\begin{itemize}[itemsep=0.3cm]
	\item \icampo
	\item \icampo
	\item \icampo
\end{itemize}

\section{Tendências do Setor}

\instrucao{Inovações, mudanças de comportamento, novas regulamentações.}

\clearpage

% -------------------------------------------------------
%  4. PRODUTOS E SERVIÇOS
% -------------------------------------------------------
\chapter{Produtos e Serviços}

\section{Descrição do Portfólio}

\centering
\begin{tabularx}{\textwidth}{X X X}
\toprule
\rowcolor{azul-claro}
\textbf{Produto / Serviço} & \textbf{Características} & \textbf{Benefícios} \\
\midrule
\tcampo{} & \tcampo{} & \tcampo{} \\
\tcampo{} & \tcampo{} & \tcampo{} \\
\tcampo{} & \tcampo{} & \tcampo{} \\
\bottomrule
\end{tabularx}
\label{tab:portfolio}

\section{Diferenciais e Inovação}

\instrucao{O que torna seu produto/serviço diferente da concorrência? Há inovação tecnológica, de processo ou de modelo de negócios?}

\section{Utilidade e Apelo ao Cliente}

\instrucao{Qual a utilidade prática? Qual o apelo emocional ou racional para o cliente?}

\section{Ciclo de Vida dos Produtos}

\centering
\begin{tabularx}{\textwidth}{X c}
\toprule
\rowcolor{azul-claro}
\textbf{Produto} & \textbf{Estágio} \\
\midrule
\tcampo{} & \tcampo{} \\
\tcampo{} & \tcampo{} \\
\tcampo{} & \tcampo{} \\
\bottomrule
\end{tabularx}

\section{Tecnologia e Propriedade Intelectual}

\instrucao{Tecnologia empregada, registro de marca, patentes, direitos autorais, P\&D.}

\section{Segurança e Conformidade}

\instrucao{Como a empresa garante a proteção de dados dos clientes e conformidade com legislação aplicável (ex.: LGPD, PCI-DSS)? Quais mecanismos técnicos e organizacionais são utilizados?}

\section{Mecanismos de Segurança Adicionais}

\instrucao{Funcionalidades extras de segurança: botão de pânico, validação de identidade, rastreamento, auditoria, etc.}

\clearpage

% -------------------------------------------------------
%  5. PLANO DE MARKETING
% -------------------------------------------------------
\chapter{Plano de Marketing}

\section{Objetivos de Marketing}

\instrucao{Metas claras e mensuráveis: market share, clientes, receita, reconhecimento de marca.}

\begin{enumerate}[itemsep=0.3cm]
	\item \icampo
	\item \icampo
	\item \icampo
\end{enumerate}

\section{Posicionamento de Mercado}

\instrucao{Como a empresa quer ser percebida pelo cliente? (Ex.: premium, custo-benefício, nicho, conveniência.)}

\section{Estratégia de Marketing Mix (4Ps)}

\subsection{Produto}

\instrucao{Embalagem, marca, qualidade, garantia, serviços agregados, design.}

\subsection{Preço}

\centering
\begin{tabularx}{\textwidth}{X r}
\toprule
\rowcolor{azul-claro}
\textbf{Produto / Serviço} & \textbf{Preço Sugerido} \\
\midrule
\tcampo{} & R\$ \tcampo{} \\
\tcampo{} & R\$ \tcampo{} \\
\tcampo{} & R\$ \tcampo{} \\
\bottomrule
\end{tabularx}
\label{tab:tabela-precos}

\subsection{Praça (Distribuição)}

\instrucao{Lojas físicas, e-commerce, marketplaces, representantes, delivery, canais digitais.}

\subsection{Promoção}

\instrucao{Estratégias de divulgação: propaganda, marketing digital, redes sociais, relações públicas, promoção de vendas, endomarketing.}

\begin{itemize}[itemsep=0.2cm]
	\item \icampo
	\item \icampo
	\item \icampo
	\item \icampo
\end{itemize}

\section{Cronograma de Marketing}

\centering
\begin{tabularx}{\textwidth}{X X r}
\toprule
\rowcolor{azul-claro}
\textbf{Período} & \textbf{Ação} & \textbf{Custo} \\
\midrule
Mês 1       & \tcampo{} & R\$ \tcampo{} \\
Mês 2--3    & \tcampo{} & R\$ \tcampo{} \\
Mês 4--6    & \tcampo{} & R\$ \tcampo{} \\
Mês 7--12   & \tcampo{} & R\$ \tcampo{} \\
\bottomrule
\end{tabularx}

\section{Projeção de Vendas}

\centering
\begin{tabularx}{\textwidth}{X r r r}
\toprule
\rowcolor{azul-claro}
\textbf{Indicador} & \textbf{Ano 1} & \textbf{Ano 2} & \textbf{Ano 3} \\
\midrule
Clientes / mês     & \tcampo{} & \tcampo{} & \tcampo{} \\
Ticket médio       & \tcampo{} & \tcampo{} & \tcampo{} \\
Faturamento mensal & \tcampo{} & \tcampo{} & \tcampo{} \\
Market-share       & \tcampo{} & \tcampo{} & \tcampo{} \\
\bottomrule
\end{tabularx}
\label{tab:projecao-vendas}

\clearpage

% -------------------------------------------------------
%  6. PLANO OPERACIONAL
% -------------------------------------------------------
\chapter{Plano Operacional}

\section{Estrutura Organizacional}

\instrucao{Organograma, cargos, funções e responsabilidades.}

\centering
\begin{tabularx}{\textwidth}{X X c}
\toprule
\rowcolor{azul-claro}
\textbf{Cargo} & \textbf{Função} & \textbf{Salário (R\$)} \\
\midrule
\tcampo{} & \tcampo{} & \tcampo{} \\
\tcampo{} & \tcampo{} & \tcampo{} \\
\tcampo{} & \tcampo{} & \tcampo{} \\
\tcampo{} & \tcampo{} & \tcampo{} \\
\bottomrule
\end{tabularx}

\section{Processos Operacionais}

\instrucao{Fluxograma dos processos de produção ou prestação de serviços.}

\begin{enumerate}[itemsep=0.3cm]
	\item \icampo
	\item \icampo
	\item \icampo
	\item \icampo
	\item \icampo
\end{enumerate}

\section{Infraestrutura e Capacidade Produtiva}

\centering
\begin{tabularx}{\textwidth}{X c c}
\toprule
\rowcolor{azul-claro}
\textbf{Item} & \textbf{Qtde} & \textbf{Valor Unit.} \\
\midrule
\tcampo{} & \tcampo{} & R\$ \tcampo{} \\
\tcampo{} & \tcampo{} & R\$ \tcampo{} \\
\tcampo{} & \tcampo{} & R\$ \tcampo{} \\
\tcampo{} & \tcampo{} & R\$ \tcampo{} \\
\bottomrule
\end{tabularx}

\vspace{0.3cm}
\textbf{Capacidade produtiva:} \icampo

\section{Fornecedores}

\centering
\begin{tabularx}{\textwidth}{X X r}
\toprule
\rowcolor{azul-claro}
\textbf{Fornecedor} & \textbf{Produto / Serviço} & \textbf{Custo Mensal} \\
\midrule
\tcampo{} & \tcampo{} & R\$ \tcampo{} \\
\tcampo{} & \tcampo{} & R\$ \tcampo{} \\
\tcampo{} & \tcampo{} & R\$ \tcampo{} \\
\bottomrule
\end{tabularx}

\section{Cronograma de Implementação}

\begin{tabularx}{\textwidth}{X X X}
\toprule
\rowcolor{azul-claro}
\textbf{Atividade} & \textbf{Prazo} & \textbf{Responsável} \\
\midrule
\tcampo{} & \tcampo{} & \tcampo{} \\
\tcampo{} & \tcampo{} & \tcampo{} \\
\tcampo{} & \tcampo{} & \tcampo{} \\
\tcampo{} & \tcampo{} & \tcampo{} \\
\bottomrule
\end{tabularx}

\section{Infraestrutura MVP}

\instrucao{Descreva a infraestrutura mínima necessária para o funcionamento do MVP (produto mínimo viável): servidores, ferramentas, serviços cloud, domínio, certificados, etc.}

\clearpage

% -------------------------------------------------------
%  7. EQUIPE DE GESTÃO
% -------------------------------------------------------
\chapter{Equipe de Gestão}

\section{Perfil dos Sócios e Executivos}

\instrucao{Nome, formação, experiência relevante, papel na empresa.}

\centering
\begin{tabularx}{\textwidth}{X X X}
\toprule
\rowcolor{azul-claro}
\textbf{Nome} & \textbf{Cargo} & \textbf{Experiência} \\
\midrule
\tcampo{} & \tcampo{} & \tcampo{} \\
\tcampo{} & \tcampo{} & \tcampo{} \\
\tcampo{} & \tcampo{} & \tcampo{} \\
\bottomrule
\end{tabularx}

\section{Conselho Consultivo e Mentores}

\instrucao{Consultores, mentores, conselheiros que apoiam o negócio.}

\section{Previsão de Crescimento de RH}

\centering
\begin{tabularx}{\textwidth}{X c c c}
\toprule
\rowcolor{azul-claro}
\textbf{Cargo} & \textbf{Ano 1} & \textbf{Ano 2} & \textbf{Ano 3} \\
\midrule
\tcampo{} & \tcampo{} & \tcampo{} & \tcampo{} \\
\tcampo{} & \tcampo{} & \tcampo{} & \tcampo{} \\
\tcampo{} & \tcampo{} & \tcampo{} & \tcampo{} \\
\textbf{Total Funcionários} & \tcampo{} & \tcampo{} & \tcampo{} \\
\bottomrule
\end{tabularx}

\clearpage

% -------------------------------------------------------
%  8. PLANO FINANCEIRO
% -------------------------------------------------------
\chapter{Plano Financeiro}

\section{Investimento Inicial}

\centering
\begin{tabularx}{\textwidth}{X r}
\toprule
\rowcolor{azul-claro}
\textbf{Item} & \textbf{Valor (R\$)} \\
\midrule
Equipamentos e Máquinas       & \tcampo{} \\
Móveis e Utensílios           & \tcampo{} \\
Reforma / Obras / Instalação  & \tcampo{} \\
Estoque Inicial               & \tcampo{} \\
Capital de Giro               & \tcampo{} \\
Despesas Pré-Operacionais     & \tcampo{} \\
\midrule
\textbf{TOTAL}                & \tcampo{} \\
\bottomrule
\end{tabularx}
\label{tab:investimento-inicial}

\section{Fontes de Recursos}

\centering
\begin{tabularx}{\textwidth}{X r X}
\toprule
\rowcolor{azul-claro}
\textbf{Fonte} & \textbf{Valor (R\$)} & \textbf{Participação} \\
\midrule
Recursos Próprios                  & \tcampo{} & \tcampo{} \\
Investidor Anjo / Venture Capital  & \tcampo{} & \tcampo{} \\
Financiamento Bancário             & \tcampo{} & \tcampo{} \\
Subvenção / Edital                 & \tcampo{} & \tcampo{} \\
\midrule
\textbf{TOTAL}                     & \tcampo{} & 100\,\% \\
\bottomrule
\end{tabularx}

\section{Premissas Financeiras}

\instrucao{Crescimento, inflação, impostos, prazos de recebimento e pagamento.}

\centering
\begin{tabularx}{\textwidth}{X X}
\toprule
\rowcolor{azul-claro}
\textbf{Premissa} & \textbf{Valor} \\
\midrule
Crescimento anual de vendas           & \tcampo{} \\
Inflação projetada                    & \tcampo{} \\
Impostos sobre faturamento            & \tcampo{} \\
Prazo médio de recebimento            & \tcampo{} \\
Prazo médio de pagamento              & \tcampo{} \\
TMA (Taxa Mínima de Atratividade)     & \tcampo{} \\
\bottomrule
\end{tabularx}

\section{Projeção de Custos e Despesas}

\subsection{Pré-Lançamento (Ano 1)}

\instrucao{Período de desenvolvimento e validação, sem receita comercial. Inclua apenas custos essenciais.}

\centering
\begin{tabularx}{\textwidth}{X r r}
\toprule
\rowcolor{azul-claro}
\textbf{Item} & \textbf{Mensal (R\$)} & \textbf{Anual (R\$)} \\
\midrule
\multicolumn{3}{@{}l}{\textbf{Custos Caixa}} \\
\quad \tcampo{}                        & \tcampo{} & \tcampo{} \\
\quad \tcampo{}                        & \tcampo{} & \tcampo{} \\
\quad \tcampo{}                        & \tcampo{} & \tcampo{} \\
\midrule
\textbf{Subtotal Caixa}                            & \tcampo{} & \tcampo{} \\
\midrule
\multicolumn{3}{@{}l}{\textbf{Custos Não-Caixa}} \\
\quad Sweat Equity (se aplicável)      & \tcampo{} & \tcampo{} \\
\midrule
\textbf{TOTAL}                                     & \tcampo{} & \tcampo{} \\
\bottomrule
\end{tabularx}
\label{tab:custos-pre-lancamento}

\subsection{Pós-Lançamento (Ano 2)}

\instrucao{Com o início da operação comercial, inclua pró-labore, marketing, infraestrutura escalada.}

\centering
\begin{tabularx}{\textwidth}{X r r}
\toprule
\rowcolor{azul-claro}
\textbf{Item} & \textbf{Mensal (R\$)} & \textbf{Anual (R\$)} \\
\midrule
\multicolumn{3}{@{}l}{\textbf{Custos Fixos}} \\
\quad \tcampo{}                        & \tcampo{} & \tcampo{} \\
\quad \tcampo{}                        & \tcampo{} & \tcampo{} \\
\quad \tcampo{}                        & \tcampo{} & \tcampo{} \\
\quad \tcampo{}                        & \tcampo{} & \tcampo{} \\
\quad \tcampo{}                        & \tcampo{} & \tcampo{} \\
\midrule
\textbf{TOTAL}                                     & \tcampo{} & \tcampo{} \\
\bottomrule
\end{tabularx}
\label{tab:custos-pos-lancamento-ano2}

\subsection{Pós-Lançamento (Ano 3)}

\instrucao{Projeção com escala ampliada.}

\centering
\begin{tabularx}{\textwidth}{X r r}
\toprule
\rowcolor{azul-claro}
\textbf{Item} & \textbf{Mensal (R\$)} & \textbf{Anual (R\$)} \\
\midrule
\multicolumn{3}{@{}l}{\textbf{Custos Fixos}} \\
\quad \tcampo{}                        & \tcampo{} & \tcampo{} \\
\quad \tcampo{}                        & \tcampo{} & \tcampo{} \\
\quad \tcampo{}                        & \tcampo{} & \tcampo{} \\
\quad \tcampo{}                        & \tcampo{} & \tcampo{} \\
\quad \tcampo{}                        & \tcampo{} & \tcampo{} \\
\midrule
\textbf{TOTAL}                                     & \tcampo{} & \tcampo{} \\
\bottomrule
\end{tabularx}
\label{tab:custos-pos-lancamento-ano3}

\section{Projeção de Faturamento}

\centering
\begin{tabularx}{\textwidth}{X r r r}
\toprule
\rowcolor{azul-claro}
\textbf{Produto / Serviço} & \textbf{Preço Unit.} & \textbf{Qtde/mês} & \textbf{Receita Mensal} \\
\midrule
\tcampo{} & \tcampo{} & \tcampo{} & \tcampo{} \\
\tcampo{} & \tcampo{} & \tcampo{} & \tcampo{} \\
\tcampo{} & \tcampo{} & \tcampo{} & \tcampo{} \\
\midrule
\textbf{Total} & & & \tcampo{} \\
\bottomrule
\end{tabularx}
\label{tab:projecao-faturamento}

\section{Demonstrativo de Resultados (DRE) Projetado}

\instrucao{Evolução dos resultados financeiros para 3 anos.}

\centering
\begin{tabularx}{\textwidth}{X r r r}
\toprule
\rowcolor{azul-claro}
\textbf{Item} & \textbf{Ano 1} & \textbf{Ano 2} & \textbf{Ano 3} \\
\midrule
1. Receita Bruta                    & \tcampo{} & \tcampo{} & \tcampo{} \\
2. (-{}) Deduções (Impostos)        & \tcampo{} & \tcampo{} & \tcampo{} \\
\midrule
3. = Receita Líquida               & \tcampo{} & \tcampo{} & \tcampo{} \\
4. (-{}) Custos Operacionais        & \tcampo{} & \tcampo{} & \tcampo{} \\
\midrule
5. = Lucro Bruto                   & \tcampo{} & \tcampo{} & \tcampo{} \\
6. (-{}) Despesas Operacionais      & \tcampo{} & \tcampo{} & \tcampo{} \\
\midrule
7. = Lucro Líquido                 & \tcampo{} & \tcampo{} & \tcampo{} \\
8. Margem Líquida                  & \tcampo{} & \tcampo{} & \tcampo{} \\
\bottomrule
\end{tabularx}
\label{tab:dre-projetado}

\section{Fluxo de Caixa Projetado}

\instrucao{Entradas e saídas de caixa previstas para os primeiros 12 meses.}

\begin{table}[htbp]
\caption{Fluxo de Caixa -- Primeiro Ano}
\label{tab:fluxo-caixa}
\footnotesize
\begin{tabularx}{\textwidth}{X *{12}{c}}
\toprule
\rowcolor{azul-claro}
\textbf{Item} & \textbf{Jan} & \textbf{Fev} & \textbf{Mar} & \textbf{Abr} & \textbf{Mai} & \textbf{Jun} & \textbf{Jul} & \textbf{Ago} & \textbf{Set} & \textbf{Out} & \textbf{Nov} & \textbf{Dez} \\
\midrule
Saldo Inicial  & & & & & & & & & & & & \\
(+) Entradas   & & & & & & & & & & & & \\
(-) Saídas     & & & & & & & & & & & & \\
\midrule
\textbf{Saldo Final} & & & & & & & & & & & & \\
\bottomrule
\end{tabularx}
\legend{Fonte: elaborado pelo autor.}
\end{table}

\section{Indicadores de Viabilidade}

\subsection{Ponto de Equilíbrio (Breakeven)}

\[
\text{Ponto de Equilíbrio} = \frac{\text{Custos Fixos Totais}}{\text{Margem de Contribuição Unitária}}
= \frac{\text{R\$ }\,\tcampo{}}{\text{R\$ }\,\tcampo{}}
= \tcampo{}\ \text{unidades/mês}
\]

\subsection{Payback}

\centering
\begin{tabularx}{\textwidth}{X r}
\toprule
\rowcolor{azul-claro}
\textbf{Item} & \textbf{Valor} \\
\midrule
Investimento Inicial            & R\$ \tcampo{} \\
Fluxo de Caixa Médio Mensal     & R\$ \tcampo{} \\
\textbf{Payback}                & \tcampo{} \\
\bottomrule
\end{tabularx}

\subsection{Valor Presente Líquido (VPL)}

\[
VPL = \sum_{t=1}^{n} \frac{FC_t}{(1 + i)^t} - I_0
= \tcampo{}
\]

\instrucao{VPL positivo indica viabilidade econômica.}

\subsection{Taxa Interna de Retorno (TIR)}

\[
\text{TIR} = \tcampo{}\ \qquad
\text{TMA} = \tcampo{}
\]

\instrucao{Se TIR $>$ TMA, o investimento é viável.}

\subsection{Resumo dos Indicadores}

\centering
\begin{tabularx}{\textwidth}{X r}
\toprule
\rowcolor{azul-claro}
\textbf{Indicador} & \textbf{Resultado} \\
\midrule
Investimento Inicial      & R\$ \tcampo{} \\
Faturamento Ano 1         & R\$ \tcampo{} \\
Lucro Líquido Ano 1       & R\$ \tcampo{} \\
Margem Líquida            & \tcampo{} \\
Payback                   & \tcampo{} \\
VPL                       & R\$ \tcampo{} \\
TIR                       & \tcampo{} \\
Ponto de Equilíbrio       & \tcampo{} \\
Índice de Lucratividade   & \tcampo{} \\
\bottomrule
\end{tabularx}
\label{tab:indicadores-viabilidade}

\clearpage

% -------------------------------------------------------
%  9. ANÁLISE DE CENÁRIOS E RISCOS
% -------------------------------------------------------
\chapter{Análise de Cenários e Riscos}

\section{Projeção de Cenários}

\instrucao{Avalie o negócio em três cenários: pessimista, realista e otimista.}

\centering
\begin{tabularx}{\textwidth}{X c c c}
\toprule
\rowcolor{azul-claro}
\textbf{Variável} & \textbf{Pessimista} & \textbf{Realista} & \textbf{Otimista} \\
\midrule
Clientes / mês      & \tcampo{} & \tcampo{} & \tcampo{} \\
Ticket médio        & \tcampo{} & \tcampo{} & \tcampo{} \\
Faturamento mensal  & \tcampo{} & \tcampo{} & \tcampo{} \\
Margem líquida      & \tcampo{} & \tcampo{} & \tcampo{} \\
Payback             & \tcampo{} & \tcampo{} & \tcampo{} \\
\bottomrule
\end{tabularx}

\section{Identificação e Classificação de Riscos}

\instrucao{Classifique os riscos por probabilidade e impacto para priorizar a mitigação.}

\centering
\begin{tabularx}{\textwidth}{X c c X}
\toprule
\rowcolor{azul-claro}
\textbf{Risco} & \textbf{Prob.} & \textbf{Impacto} & \textbf{Estratégia de Mitigação} \\
\midrule
\tcampo{} & \tcampo{} & \tcampo{} & \tcampo{} \\
\tcampo{} & \tcampo{} & \tcampo{} & \tcampo{} \\
\tcampo{} & \tcampo{} & \tcampo{} & \tcampo{} \\
\tcampo{} & \tcampo{} & \tcampo{} & \tcampo{} \\
\bottomrule
\end{tabularx}

\section{Plano de Contingência}

\instrucao{Ações pré-definidas caso um risco crítico se concretize.}

\clearpage

% -------------------------------------------------------
%  10. ESTRATÉGIA DE CRESCIMENTO
% -------------------------------------------------------
\chapter{Estratégia de Crescimento}

\section{Análise SWOT}

\centering
\begin{tabularx}{\textwidth}{X X}
\toprule
\rowcolor{azul-claro}
\textbf{Forças (Strengths)} & \textbf{Fraquezas (Weaknesses)} \\
\midrule
1. \tcampo{} & 1. \tcampo{} \\
2. \tcampo{} & 2. \tcampo{} \\
3. \tcampo{} & 3. \tcampo{} \\
\addlinespace
\rowcolor{azul-claro}
\textbf{Oportunidades (Opportunities)} & \textbf{Ameaças (Threats)} \\
\midrule
1. \tcampo{} & 1. \tcampo{} \\
2. \tcampo{} & 2. \tcampo{} \\
3. \tcampo{} & 3. \tcampo{} \\
\bottomrule
\end{tabularx}

\section{Objetivos e Metas SMART}

\instrucao{Específicas, Mensuráveis, Atingíveis, Relevantes, com Prazo.}

\centering
\begin{tabularx}{\textwidth}{X X c}
\toprule
\rowcolor{azul-claro}
\textbf{Objetivo} & \textbf{Meta} & \textbf{Prazo} \\
\midrule
\tcampo{} & \tcampo{} & \tcampo{} \\
\tcampo{} & \tcampo{} & \tcampo{} \\
\tcampo{} & \tcampo{} & \tcampo{} \\
\bottomrule
\end{tabularx}

\section{Estratégia de Entrada e Crescimento}

\instrucao{Matriz Produto-Mercado (Ansoff): penetração, desenvolvimento de produto, expansão de mercado ou diversificação?}

\centering
\begin{tabularx}{\textwidth}{X X}
\toprule
\rowcolor{azul-claro}
\textbf{Tipo} & \textbf{Descrição} \\
\midrule
Penetração de Mercado       & \tcampo{} \\
Desenvolvimento de Produto  & \tcampo{} \\
Expansão de Mercado         & \tcampo{} \\
Diversificação              & \tcampo{} \\
\bottomrule
\end{tabularx}

\section{Riscos Críticos do Negócio}

\instrucao{Fatores que podem inviabilizar o negócio e como mitigá-los.}

\clearpage

% -------------------------------------------------------
%  11. CONSIDERAÇÕES FINAIS
% -------------------------------------------------------
\chapter{Considerações Finais}

\instrucao{Reflexões finais sobre o potencial do negócio, impacto esperado e visão de futuro.}

\vspace{0.5cm}

\noindent\textbf{Conclusão:}

\vspace{2cm}

\begin{table}[htbp]
\centering
\caption{Cenário Realista -- Projeção consolidada}
\label{tab:cenario-realista}
\centering
\begin{tabularx}{\textwidth}{X r}
\toprule
\rowcolor{azul-claro}
\textbf{Indicador} & \textbf{Resultado} \\
\midrule
Investimento Inicial & R\$ \tcampo{} \\
Faturamento Ano 1    & R\$ \tcampo{} \\
Faturamento Ano 2    & R\$ \tcampo{} \\
Faturamento Ano 3    & R\$ \tcampo{} \\
Lucro Líquido Ano 1  & R\$ \tcampo{} \\
Lucro Líquido Ano 2  & R\$ \tcampo{} \\
Lucro Líquido Ano 3  & R\$ \tcampo{} \\
Margem Líquida (Ano 3) & \tcampo{} \\
Payback              & \tcampo{} \\
VPL                  & R\$ \tcampo{} \\
TIR                  & \tcampo{} \\
Ponto de Equilíbrio  & \tcampo{} \\
\bottomrule
\end{tabularx}
\end{table}

\noindent\textbf{Próximos passos:}
\begin{enumerate}[itemsep=0.3cm]
	\item \icampo
	\item \icampo
	\item \icampo
\end{enumerate}

\vspace{1.5cm}

\begin{flushright}
	[Cidade], \underline{\hspace{4cm}} de \underline{\hspace{3cm}} de \underline{\hspace{2cm}}. \\[1.5cm]
	\rule{7cm}{0.4pt} \\[0.3cm]
	\textbf{\campo[7cm]{}} \\
	\textit{Empreendedor / Responsável}
\end{flushright}

\clearpage

% -------------------------------------------------------
%  12. REFERÊNCIAS
% -------------------------------------------------------
\chapter{Referências}

\instrucao{Liste aqui todas as fontes consultadas para elaboração do plano: dados de mercado, relatórios setoriais, sites, livros.}

\begin{enumerate}[itemsep=0.3cm]
	\item \icampo
	\item \icampo
	\item \icampo
	\item \icampo
	\item \icampo
\end{enumerate}

\clearpage

% -------------------------------------------------------
%  13. ANEXOS
% -------------------------------------------------------
\chapter{Anexos}

\instrucao{Documentos complementares: CVs, pesquisas, contratos, fotos, plantas baixas, layout, etc.}

\begin{enumerate}[itemsep=0.3cm]
	\item \icampo
	\item \icampo
	\item \icampo
	\item \icampo
	\item \icampo
\end{enumerate}

\vspace{1cm}

{
\centering
\color{cinza-instrucao}
\rule{0.6\textwidth}{0.4pt} \\[0.3cm]
\textit{--- FIM DO DOCUMENTO ---} \\[0.3cm]
\textit{O plano de negócios é um documento vivo. Revise e atualize sempre que o mercado ou o negócio mudar.}
\par
}

\end{document}
